\chapter*{General Introduction}
\addcontentsline{toc}{chapter}{General Introduction}
\markboth{General Introduction}{}

The car rental business has grown a lot because more people travel, live in cities, and want easy ways to move around. But, old-style car rental companies find it hard to manage bookings, check which cars are free, and keep customers happy. These problems show we need new, simple, and easy-to-use systems to make car rentals better and improve service. 

We introduce "RentyCar," a complete car rental system that solves these problems using new web tools and good software methods. It makes renting cars easier, helps the business work better, and gives customers and companies a nicer experience. 

RentyCar was created for several important reasons. It helps track available cars and manage bookings instantly, making the process quick and easy. The system also allows businesses to handle many rental locations smoothly and efficiently. Additionally, it ensures secure and reliable user logins to keep information safe. Finally, RentyCar provides complete tools to manage cars and offices, improving operations for both customers and companies.

This report is structured into three chapters :
\begin{itemize}
    \item \textbf{Chapter 1:} Presents the preliminary study, including how the project was planned, looking at other solutions, and the idea for RentyCar.
    
    \item \textbf{Chapter 2:} Details the conceptual study, covering physical architecture modeling, use case diagrams, class diagrams, and sequence diagrams.
    
    \item \textbf{Chapter 3:} Describes building and testing the system, including the chosen structure, tools used, and how the system was made.
\end{itemize}
Finally, we conclude this report with an analysis of the results obtained, a discussion
of the challenges encountered, and future perspectives for the evolution.

